\documentclass{jstarticle}
\usepackage[utf8]{inputenc}
\usepackage[T1]{fontenc}
\usepackage[english]{babel}
\usepackage{csquotes}
\usepackage{lmodern}
\renewcommand{\familydefault}{\rmdefault}

\usepackage{listings}
\lstset{
  basicstyle=\ttfamily\footnotesize,
  keywordstyle=\color{blue}\bfseries,
  commentstyle=\color{gray},
  stringstyle=\color{green!50!black},
  breaklines=true,
  frame=none,
  backgroundcolor=\color{gray!5},
  showstringspaces=false
}

% If you have long equations
% \documentclass[mathindent=1ex]{jstarticle}

% Which journal to submit to.
% - "ssad": Smart Systems and Devices
% - "etsd": Engineering and Technology for Sustainable Development
% - "ds-ai": Data Science and Artificial Intelligence
% Default is "ssad".
\submitTo{ssad}

% Manuscript title
\title{Adaptive Multi-Path Inference for Multi-Tenant LLM Systems with Shared Models and LoRA-based Personalization}
% Authors
\author[1]{Tran Bui Tung Linh}
\author[2]{Hoang Trung Kien}
\author[3]{Nguyen Thanh Trung}
\authornote{}{Faculty of Mathematics and Informatics, Hanoi University of Science and Technology, Ha Noi, Vietnam}

\authornote{*}{Corresponding author email: linh.tbt227057@hust.edu.vn}

% Abstract
\begin{abstract}
Enterprise LLM applications are increasingly being utilized in multi-tenant settings, specifically where heterogeneous queries from different organizations are funneled into shared infrastructures, while under a mandate to ensure both efficiency, isolation, and tenant-specific behavior. 
In such environments, single-path inference pipelines often prove insufficient, as simple and complex requests are forced through the same expensive execution process. 
This paper presents an adaptive multi-path inference architecture for multi-tenant LLM systems, capable of dynamically selecting among retrieval-based generation, tool-based processing for structured or semi-structured data, and direct generation according to query characteristics. 
The system is built around a shared base model to elevate resource efficiency, while tenant-specific Low-Rank Adaptation (LoRA) modules support lightweight tenant-level customization without replicating the full model for each tenant.
To support practical deployment, we further introduce a tenant-aware runtime with isolated knowledge stores, isolated conversational memory, and dynamic adapter loading to reduce memory overhead during inference. 
The proposed design positions retrieval and tools as execution paths within a broader serving architecture rather than as the defining identity of the system. 
Experimental evaluation in a multi-tenant setting shows that the architecture improves latency and memory efficiency over fixed single-path baselines while maintaining competitive response quality. 
These results suggest that adaptive multi-path inference, combined with shared-model serving and LoRA-based personalization, is a promising design paradigm for scalable and cost-efficient enterprise LLM systems.
\end{abstract}

% Keywords
\keyword{Multi-tenant LLM systems}
\keyword{Adaptive inference}
\keyword{LoRA personalization}
\keyword{Resource efficiency}

% Bibliography
% \usepackage[backend=biber,style=ieee]{biblatex}
\addbibresource{refs.bib}

% Other settings
% These settings are document specific.
\usepackage{tabularray}

\usepackage{xcolor}
\definecolor{customgray}{gray}{0.9}
\usepackage{tabularx}

\usepackage{listings}
\lstset{basicstyle=\ttfamily}
\let\Code\lstinline

\usepackage{prettyref}
\newrefformat{fig}{Fig.\,\ref{#1}}
\newrefformat{tab}{Table.\,\ref{#1}}
\newrefformat{sec}{Section \ref{#1}}
\newrefformat{subsec}{Section \ref{#1}}
\newrefformat{subsubsec}{Section \ref{#1}}

\begin{document}

\maketitle

\section{Introduction}

\subsection{Background and Motivation}

\noindent Large language models (LLMs) are surging in adoption, assuming the role of enterprise assistants, document intelligence, customer support, and internal knowledge systems. In practice, these deployments rarely serve a single homogeneous workload. Instead, they must process heterogeneous requests ranging from general conversational questions to document-grounded queries and structured data operations. This heterogeneity becomes more challenging in multi-tenant environments, where multiple companies and organizations share the same serving infrastructure but still expect tenant-specific behavior, data isolation, and predictable performance.

Many practical deployments still rely solely on a fixed inference pipeline for all requests. Although such a design is simple to implement, it is inefficient for enterprise workloads as different query types require different execution flows. Routing every request through the same expensive retrieval or generation process can lower response rates, waste compute resources, and blur the system-level trade-off between answer quality and serving efficiency.

\subsection{Research Gaps in Multi-tenant LLM Serving}

The first gap is an \textit{efficiency gap}. Existing one-path-fits-all pipelines do not distinguish well between simple queries, knowledge-grounded queries, and structured data queries. As a result, lightweight requests may still trigger retrieval, long-context prompting, or unnecessary model calls.

The second gap is a \textit{personalization gap}. Multi-tenant applications often require tenant-specific language style, policies, and domain preferences. A straightforward solution is to maintain separate models for different tenants, but this quickly becomes impractical because full-model duplication increases memory cost and complicates system maintenance.

The third gap is a \textit{multi-tenant systems gap}. A realistic serving system must jointly support tenant isolation, scalable resource sharing, and personalization under the same runtime. Many existing approaches optimize only one dimension, such as retrieval quality, adapter efficiency, or serving throughput, without addressing the interaction among routing, shared serving, and tenant-aware isolation in a unified architecture.

\subsection{Adaptive Multi-Path Serving Architecture}

To address these gaps, this paper frames the problem as adaptive multi-path serving for multi-tenant LLM systems. Rather than treating retrieval-augmented generation as the default identity of the system, we treat retrieval as one execution path among several possible paths. The proposed architecture selects between tool-based processing, retrieval-based generation, and direct generation based on query characteristics and tenant context.

The system adopts a shared base model as the common serving backbone, which improves scalability and reduces memory usage compared with per-tenant model replication. To preserve tenant-specific behavior under shared serving constraints, we incorporate Parameter-Efficient Fine-Tuning (PEFT) through Low-Rank Adaptation (LoRA). In this design, LoRA is not merely an auxiliary personalization module; it is the practical mechanism that resolves the trade-off between shared-model efficiency and tenant-level customization.

The runtime is tenant-aware by construction. Each tenant is associated with isolated knowledge storage, isolated conversational memory, and an optional personalization adapter. Tenant-specific LoRA adapters are loaded dynamically during inference so that the system can support customization without keeping all adapters resident in memory at the same time.

\subsection{Scope and Contributions}

The paper is positioned as a systems-and-application study on multi-tenant LLM serving. Its central claim is not that retrieval, tools, or LoRA are individually new, but that their coordinated integration yields a practical serving architecture for heterogeneous enterprise workloads.

\subsection{Contributions}

The main contributions of this work are summarized as follows:

\begin{itemize}
    \item We propose a unified adaptive multi-path serving architecture for heterogeneous multi-tenant LLM workloads, where queries are routed among tool, retrieval, and direct-generation paths.
    \item We integrate shared-model serving with tenant-specific LoRA personalization to balance resource efficiency and customization without full-model duplication.
    \item We introduce a tenant-aware runtime that combines isolated knowledge storage, isolated memory, and dynamic adapter loading for scalable multi-tenant deployment.
    \item We evaluate the architecture using system-oriented metrics, including latency, memory usage, and response quality, to compare adaptive serving against fixed single-path baselines.
\end{itemize}

\section{System Architecture}
\subsection{Overall Structure}

The proposed system is designed as a modular multi-tenant serving architecture that shares the core model across tenants while preserving tenant-specific data boundaries. Its main components are the data ingestion layer, the tenant-aware storage and retrieval layer, the adaptive inference engine, and the API serving layer.

At the system level, the architecture is not organized around a single fixed RAG pipeline. Instead, it uses adaptive multi-path inference, where retrieval is one possible path inside a broader serving framework. A shared base language model is deployed as the main inference engine and serves multiple tenants through a common runtime.

Each tenant maintains an isolated knowledge base and conversation state. During query processing, the runtime first identifies the tenant context, then applies route selection, and finally invokes the most appropriate execution path under that tenant scope.

The routing layer determines whether a query should be processed through tool-based execution, retrieval-based generation, or direct generation. This design reduces unnecessary LLM calls and better matches execution cost to query complexity.

\subsection{Data Ingestion and Knowledge Base Construction}

The data ingestion layer is responsible for collecting, processing, and indexing tenant-specific data from multiple sources.

Supported data types include:
\begin{itemize}
    \item Unstructured documents (PDF, DOCX, HTML)
    \item Structured data (Excel, CSV)
    \item Web-based content (links, HTML pages)
\end{itemize}

The ingestion pipeline performs document parsing, text normalization, chunking, embedding generation, and vector storage. Each tenant's data is stored in an isolated partition to enforce separation and reduce cross-tenant leakage risk.

\subsection{Retrieval Module}

The retrieval module is responsible for fetching relevant context only when the selected route requires external grounding.

Given a user query and tenant identifier:
\begin{enumerate}
    \item The query is converted into an embedding vector
    \item Similarity search is performed within the tenant-specific index
    \item Top-k relevant chunks are returned
\end{enumerate}

To improve retrieval quality, the system can incorporate chunk filtering, re-ranking, and hybrid retrieval strategies. In the proposed architecture, retrieval is therefore a selective component of the serving process rather than the default path for every request.

\subsection{Dynamic Inference Engine}

The dynamic inference engine is the core component of the system. It combines a shared base LLM, a route-selection mechanism, and optional tenant-specific personalization modules.

\subsubsection{Shared Base Model}

A single base LLM is shared across all tenants, which reduces memory usage and simplifies deployment compared with maintaining a separate full model for each tenant.

\subsubsection{Query Routing Mechanism}

The system introduces a routing mechanism that maps incoming queries to one of several execution paths:
\begin{itemize}
    \item Tool-based queries (e.g., Excel parsing, file inspection)
    \item Retrieval-based queries (RAG)
    \item General queries handled directly by the LLM
\end{itemize}

This routing structure minimizes unnecessary model invocations and improves latency for heterogeneous workloads.

\subsubsection{Dynamic Adapter Loading}

To support tenant-level customization, the system uses PEFT-based adapters, specifically LoRA, that are dynamically loaded during inference. Instead of keeping all adapters in memory, the runtime loads the required adapter on demand and releases it after use, which lowers memory overhead while preserving personalization capability.

\subsection{API and Serving Layer}

The architecture is exposed through a RESTful API implemented with FastAPI.

Key responsibilities include:
\begin{itemize}
    \item Handling user requests
    \item Resolving tenant context
    \item Invoking the selected execution path
    \item Returning responses and associated provenance when applicable
\end{itemize}

The API layer also supports logging, monitoring, and integration with front-end applications. These functions are important for reproducibility and operational visibility in multi-tenant deployments.

\section{Implementation Details}

\subsection{System Components}

The system is implemented using a modular architecture, where each component is responsible for a specific function in the query processing pipeline.

The main components include:
\begin{itemize}
    \item \textbf{Router module:} Determines the execution path for each query.
    \item \textbf{Workflow manager:} Orchestrates the end-to-end processing pipeline.
    \item \textbf{Retrieval module:} Handles tenant-scoped vector-based similarity search.
    \item \textbf{LLM service:} Manages inference using the shared base model.
    \item \textbf{Memory store:} Maintains tenant-specific data and context.
    \item \textbf{Tool module:} Processes structured data queries (e.g., Excel, CSV).
\end{itemize}

This modular design improves flexibility, scalability, and maintainability of the system.

\subsection{Multi-tenant Design}

Each tenant in the system is associated with:
\begin{itemize}
    \item A unique tenant identifier
    \item An isolated knowledge base
    \item An isolated memory context
    \item An optional LoRA personalization adapter
\end{itemize}

The system ensures strict isolation between tenants by:
\begin{itemize}
    \item Separating vector database partitions
    \item Maintaining independent memory contexts
    \item Routing queries based on tenant identifiers
\end{itemize}

At the same time, all tenants share a single base model, enabling efficient resource utilization.

\subsection{Routing Optimization}

The routing mechanism is designed to reduce unnecessary computational overhead and to better align execution cost with query requirements.

Instead of processing all queries through the LLM, the system dynamically selects the most efficient execution path:

\begin{itemize}
    \item \textbf{Tool-based processing:} For structured data queries such as Excel or CSV
    \item \textbf{Retrieval-based generation:} For document-based queries (e.g., PDF)
    \item \textbf{Direct LLM inference:} For general queries
\end{itemize}

This approach reduces inference latency and token usage by avoiding a uniform pipeline for all requests.

\subsection{Resource Optimization Strategy}

The system incorporates multiple techniques to optimize resource usage:

\begin{itemize}
    \item \textbf{Shared model architecture:} A single base LLM serves all tenants
    \item \textbf{Dynamic adapter loading:} Tenant-specific LoRA adapters are loaded on demand
    \item \textbf{Selective LLM invocation:} Retrieval and full generation are used only when justified by the route decision
\end{itemize}

These strategies reduce memory footprint and improve scalability in cloud environments.

\section{Experimental Evaluation}

\subsection{Experimental Setup}

The system is evaluated using a test set that covers heterogeneous enterprise-style queries:
\begin{itemize}
    \item Structured data queries handled by the tool path
    \item Document-grounded queries handled by the retrieval path
    \item General reasoning queries handled by direct generation
    \item Multi-tenant requests executed under isolated tenant contexts
\end{itemize}

Two configurations are compared:
\begin{itemize}
    \item Fixed single-path baselines, where queries are forced through one dominant processing pipeline
    \item The proposed adaptive system, which combines route selection, shared-model serving, and tenant-aware execution
\end{itemize}

\subsection{Evaluation Metrics}

The following metrics are used:
\begin{itemize}
    \item Latency (ms): time taken to produce a final response
    \item Memory usage (MB): RAM or GPU consumption during serving
    \item Response quality: correctness and usefulness of generated answers
    \item Path suitability: whether the selected execution path matches the query type
\end{itemize}

\subsection{Results and Analysis}

The results show that adaptive multi-path inference reduces latency by avoiding unnecessary retrieval or generation steps for lightweight queries. The shared-model design also lowers memory usage compared with configurations that replicate full models or overuse retrieval-heavy execution.

Structured queries benefit the most from routing because they can be resolved through tools without incurring the full cost of long-context LLM inference. Retrieval-heavy queries still benefit from grounding when needed, while general queries avoid unnecessary external lookup.

\subsection{Discussion}

Overall, the evaluation supports the central systems claim of the paper: adaptive route selection, shared-model serving, and lightweight personalization work together to improve efficiency while preserving multi-tenant practicality.

\section{Conclusion}

This paper presents an adaptive multi-path serving architecture for multi-tenant LLM systems with shared models and LoRA-based personalization. The proposed design addresses three coupled requirements of enterprise deployment: efficient serving under heterogeneous workloads, tenant-specific customization without full-model duplication, and tenant-aware isolation within a shared runtime.

By treating retrieval, tool use, and direct generation as coordinated execution paths instead of forcing all requests through a single pipeline, the system better aligns inference cost with query needs. The shared base model improves resource efficiency, while dynamic LoRA loading enables lightweight personalization that remains compatible with multi-tenant deployment constraints.

\subsection{Future Work}

Future work includes:
\begin{itemize}
    \item learning-based routing policies with calibrated confidence scores,
    \item stronger evaluation of personalization under tenant-specific benchmarks,
    \item larger-scale system experiments across more tenants and workloads, and
    \item deeper analysis of throughput, adapter switching overhead, and isolation guarantees in production-like settings.
\end{itemize}

% Acknowledgments and references will be added in the final submission version.
% \section*{Acknowledgments}
Use the singular heading even if you have many acknowledgments. Avoid expressions such as ``One of us (S.B.A.) would like to thank ..." Instead, write ``F. A. Author thanks ..."

It will also contain support information, including sponsor and financial support acknowledgment. For example, ``This work was supported by the ... under the Grant T2020-PC000".

% \section*{Reference}
Reference style should be strictly followed by the guidance in the examples below. References should be numbered in sequence and cited in the main text, and must be arranged in the order they are cited in the main text. When cited, the reference should appear on the line, in square brackets, inside the punctuation, with the reference number such as \cite{BookTemplate}. In case, it is necessary to cite more than one reference, write them as \cite{BookTemplate, ArticleTemplate, ReportTemplate} or \cite{BookTemplate, HanbookTemplate, ReportHanbookOnlineTemplate}. Please do not use automatic endnotes in \textit{Word}, rather, type the reference list at the end of the paper using the “References” style. All other contents remain unchanged. The first line indent of 0.75 cm applies only to the reference list.

Reference numbers are set flush left and form a column of their own, hanging out beyond the body of the reference. The reference numbers are on the line, enclosed in square brackets. In all references, the given name of the author or editor is abbreviated to the initial only and precedes the last name. Use commas around Jr., Sr., in names. Abbreviate conference titles. When referencing a patent, provide the day and the month of issue, or application. References may not include all information; please obtain and include relevant information. Do not combine references. There must be only one reference with each number. If there is a URL included with the print reference, it can be included at the end of the reference. In this case, a DOI number for the full text is recommended to use.

References in other languages should be translated into English and follow the JST following reference format and the title in original languages need to be provided, Example: N. T. Nguyen, T. H. Tran, and L. Q. Pham, Impact of the Fourth Industrial Revolution on Vietnam's textile and garment industry], Journal of Science and Technology, vol. 58, no. 3, pp. 45--56, 2023. [In Vietnames]: Ảnh hưởng của Cách mạng công nghiệp lần thứ tư đến ngành dệt may Việt Nam.

Other than books, capitalize only the first word in a paper title, except for proper nouns and element symbols. For papers published in translation journals, please give the English citation first, followed by the original foreign-language citation. See the end of this document for formats and examples of common references.

\begin{itemize}

  \item \textit{\textbf{Basic format for books:}} \cite{BookTemplate}
  \begin{lstlisting}[language=TeX]
@book{BookTemplate,
    author = {Author, J. K.},
    title = {Title of His Published Book},
    chapter = {x},
    section = {x},
    chaptertitle = {Title of chapter in the book},
    edition = {n},
    location = {City of Publisher, Country},
    publisher = {Abbrev. of Publisher},
    address = {City of Publisher, Country},
    date = {1970},
    pages = {xx--yy},
    pagetotal = {ppp},
    keywords = {template},
}
  \end{lstlisting}

  \item \textit{\textbf{Basic format for periodicals:}} \cite{ArticleTemplate}
  \begin{lstlisting}[language=TeX]
@article{ArticleTemplate,
    keywords = {template},
    author = {Author, J. K.},
    title = {Name of paper},
    journal = {Abbrev. title of periodical},
    volume = {11},
    number = {2},
    pages = {123-456},
    month = {1},
    year = {2024},
    doi = {https://doi.org/10.1109. xx.123456},
}
  \end{lstlisting}

  \item \textit{\textbf{Basic format for reports:}} \cite{ReportTemplate}
  \begin{lstlisting}[language=TeX]
@report{ReportTemplate,
    author = {Author, J. K.},
    title = {Title of report},
    institution = {Name of Co.},
    address = {City of Co., Abbrev. State, Country},
    year = {2024},
}
  \end{lstlisting}

\newpage
  \item \textit{\textbf{Basic format for handbooks:}} \cite{HanbookTemplate}
  \begin{lstlisting}[language=TeX]
@book{HanbookTemplate,
    author    = {Author Name(s)},
    title     = {Title of the Handbook},
    edition   = {Edition (if applicable)},
    year      = {Year of Publication},
    publisher = {Publisher Name},
    address   = {Publisher Location},
    isbn      = {ISBN Number (if available)}
}
  \end{lstlisting}

  \item \textit{\textbf{Basic format for reports and handbooks (when available online):}} \cite{ReportHanbookOnlineTemplate}
  \begin{lstlisting}[language=TeX]
@techreport{ReportHanbookOnlineTemplate,
    author    = {Author Name(s)},
    title     = {Title of the Report/Handbook},
    institution = {Institution/Organization Name},
    year      = {Year of Publication},
    number    = {Report Number (if applicable)},
    type      = {Report Type (e.g., Technical Report)},
    url       = {URL},
    note      = {Accessed: Month Day, Year}
}
  \end{lstlisting}

  \item \textit{\textbf{Basic format for computer programs and electronic documents (when available online):}} \cite{MISCTemplate}
  \begin{lstlisting}[language=TeX]
@misc{MISCTemplate,
    author    = {Author Name(s) or Organization},
    title     = {Title of the Software or Document},
    year      = {Year of Release or Update},
    url       = {URL},
    note      = {Version number, Accessed: Month Day, Year}
}
  \end{lstlisting}

\newpage
  \item \textit{\textbf{Example for papers presented at conferences (unpublished):}} \cite{InproceedingTemplate}
  \begin{lstlisting}[language=TeX]
@inproceedings{InproceedingTemplate,
    keywords = {template},
    author = {Author, J. K.},
    title = {Title of paper},
    booktitle = {Abbreviated Name of Conf.},
    location = {City of Conf., Abbrev. State (if given), Country},
    date = {1970},
    pages = {xx-xx},
}
  \end{lstlisting}

  \item \textit{\textbf{Basic format for patents:}} \cite{PatentTemplate}
  \begin{lstlisting}[language=TeX]
@patent{PatentTemplate,
    author = {Author, J. K.},
    title = {Title of patent},
    type = {U. S. Patent},
    number = {xx},
    date = {1970-12-31},
}
  \end{lstlisting}

  \item \textit{\textbf{Basic format for theses (M.S.) and dissertations (Ph.D.):}} \cite{MasterThesisTemplate}, \cite{PhDThesisTemplate}
  \begin{lstlisting}[language=TeX]
@mastersthesis{MasterThesisTemplate,
    author = {Author, J. K.},
    title = {Title of thesis},
    institution = {Abbrev. Dept, Abbrev. Univ.},
    school = {Abbrev. Dept, Abbrev. Univ.},
    location = {City of Univ., Abbrev. State},
    date = {1970},
    sortkey = {1},
}
  \end{lstlisting}

  \begin{lstlisting}[language=TeX]
@phdthesis{PhDThesisTemplate,
    author = {Author, J. K.},
    title = {Title of dissertation},
    institution = {Abbrev. Dept, Abbrev. Univ.},
    school = {Abbrev. Dept, Abbrev. Univ.},
    location = {City of Univ., Abbrev. State},
    date = {1970},
    sortkey = {2},
}
  \end{lstlisting}

  \item \textit{\textbf{Basic formats for standards:}} \cite{IEEE:Standard308}
  \begin{lstlisting}[language=TeX]
@standard{IEEE:Standard308,
    author = {IEEE},
    title = {IEEE Criteria for Class IE Electric Systems},
    number = {IEEE Standard 308},
    year = {1969},
}
  \end{lstlisting}

\end{itemize}

% \printbibliography[heading=none]
% \printbibliography
\end{document}
